\section{Relationships and Mobilizing Help}

\subsection{Beyond the Vanity Rolodex}

The phrase \textit{I have a network} usually conceals more than it reveals. It might mean hundreds of superficial connections, 10 trusted collaborators, 3 purchasing decision-makers, or 1 person willing to stake their reputation on you. These are not interchangeable.

Research on social networks demonstrates that relationships create access to valuable information and opportunities beyond one's immediate circle. A large randomized field experiment on LinkedIn found causal effects of network tie strength on job mobility, with variation by industry and by how tie strength was measured.\footnote{Karthik Rajkumar, Guillaume Saint-Jacques, Iavor Bojinov, Erik Brynjolfsson, and Sinan Aral, ``A causal test of the strength of weak ties,'' \textit{Science} 377, no. 6612 (2022): 1304--1310.} That finding concerns employment networks; it does not establish that weak ties automatically produce startup customers. It suggests a mechanism worth testing: different kinds of relationships reveal different information and open different doors.

\subsection{Audit Relationships by Function}

For each relevant person, record their operational role rather than just their name:

\begin{itemize}[leftmargin=11pt,itemsep=1.5pt,topsep=1.5pt,parsep=0pt]
  \item \textbf{Problem witness:} Experiences the operational pain firsthand.
  \item \textbf{Buyer:} Can authorize spending or influence a purchasing process.
  \item \textbf{Connector:} Can introduce you to buyers or specialists.
  \item \textbf{Builder:} Can help deliver technical or operational work.
  \item \textbf{Domain critic:} Will tell you when an idea is naive.
  \item \textbf{Credibility witness:} Can accurately verify a past result.
  \item \textbf{Capital partner:} Can contribute resources under agreed terms.
  \item \textbf{Mentor or peer:} Can challenge your reasoning on difficult decisions.
\end{itemize}

A single person can play several roles. But you should not count a person as a buyer merely because they work in the same industry.

\subsection{The Activation Test}

For each person, ask privately:

\begin{enumerate}[label=\textbf{\arabic*.},leftmargin=11pt,itemsep=1.5pt,topsep=1.5pt,parsep=0pt]
  \item Would they recognize me and reply to a respectful message?
  \item Have we exchanged something of value before?
  \item Is the relationship current, or merely historical?
  \item Is it appropriate for me to contact them about this topic?
  \item What specific, reasonable request could I make?
  \item What useful thing can I offer in return?
  \item Would involving them create a conflict of interest, privacy issue, or unwanted pressure?
\end{enumerate}

Use 3 simple statuses: \textbf{unverified}, \textbf{responsive}, \textbf{committed}. Move a relationship to \textit{committed} only when the person actually agrees to an action (a meeting, introduction, pilot, contribution, or reference).

\textbf{Counting names before activation creates an illusion of distribution; counting verified commitments creates commercial evidence.}

\subsection{The Network Map That Helps}

\begin{tcolorbox}[
  enhanced,
  colback=white,
  colframe=BayesBorder!80!white,
  boxrule=0.4pt,
  arc=2pt,
  left=3pt,right=3pt,top=3pt,bottom=3pt,
  before skip=3pt,after skip=5pt
]
\scriptsize
\setlength{\tabcolsep}{2pt}
\renewcommand{\arraystretch}{1.18}
\begin{tabularx}{\linewidth}{@{} >{\bfseries\color{BayesInk}}p{1.8cm} >{\centering\arraybackslash}p{1.1cm} >{\centering\arraybackslash}p{1.0cm} >{\centering\arraybackslash}p{1.0cm} >{\color{BayesSecondary}}X @{}}
\toprule
\textbf{Segment} & \textbf{Witnesses} & \textbf{Buyers} & \textbf{Intros} & \textbf{Immediate Next Action} \\
\midrule
Engineering Colleges & 8 & 1 & 2 & Map annual fiscal buying cycle \\
Mid-Size Dev Firms & 12 & 0 & 3 & Run 5 workflow interviews \\
Consumer Health & 4 & 0 & 1 & Learn category unit margins \\
\bottomrule
\end{tabularx}
\end{tcolorbox}

All numbers above are illustrative. Your actual inventory may be 0 in every column; that is useful information, not failure.

\subsection{Permission as a Separate Asset}

An address book, student list, customer database, or employer CRM is not automatically a personal acquisition channel. Access to personal information does not imply permission to reuse it. Employment agreements, privacy rules, contractual limits, and reasonable expectations may all apply.

Make permission explicit in your distribution audit: \textit{Can I contact this person? For what purpose? Through which channel? On what basis?} If the answer is uncertain, do not count that database as available distribution. Build opt-in access instead.

\textbf{Permission is a separate asset: access to personal data does not grant legal or ethical permission to commercialize it.} The best day-zero strategy is to ask well-scoped questions, do something useful, and see who voluntarily deepens the relationship.
